\section{结果与讨论}

\paragraph{多类型纳米颗粒基准数据集上的模型比较。}
为获得可靠的性能基线，我们首先在大规模多类型纳米颗粒毒性数据集上比较六类模型的表现，包括随机森林\cite{breiman2001random}、XGBoost\cite{chen2016xgboost}、ABRF\cite{gromov2022attention}、LARF\cite{gromov2022larf}、TabTransformer\cite{huang2020tabtransformer} 以及蒸馏 TabTransformer\cite{hinton2015distilling}。该基准数据集来源于 Masarkar 等人建立的公开纳米颗粒毒性数据资源 \cite{masarkar2024}，覆盖金属、金属氧化物、聚合物、碳基与无机纳米颗粒，清洗后共包含 4102 个有效样本，按照 7:3 划分为训练集 3281 条与测试集 821 条。输入特征包括浓度、暴露时间、材料类别、粒径、Zeta 电位、细胞类型与检测方法等。

回归任务结果如表~\ref{tab:benchmark_regression} 所示。TabTransformer 取得最高的决定系数，$R^2$ 为 0.768，MAE 为 10.30\%，RMSE 为 14.97\%。XGBoost 与随机森林表现接近，分别达到 $R^2=0.752$ 与 $R^2=0.749$。引入注意力机制后的 ABRF 与 LARF 未超过传统树模型，蒸馏 TabTransformer 的性能则明显下降。

\begin{table}[htbp]
  \centering
  \caption{多类型纳米颗粒基准数据集上的回归性能比较}
  \label{tab:benchmark_regression}
  \begin{tabular}{lccc}
    \hline
    模型 & $R^2$ & MAE (\%) & RMSE (\%) \\
    \hline
    TabTransformer & 0.768 & 10.30 & 14.97 \\
    XGBoost & 0.752 & 11.13 & 15.56 \\
    随机森林 & 0.749 & 10.81 & 15.71 \\
    ABRF & 0.719 & 11.87 & 16.55 \\
    LARF & 0.716 & 11.94 & 16.63 \\
    蒸馏 TabTransformer & 0.506 & 16.85 & 22.31 \\
    \hline
  \end{tabular}
\end{table}

分类任务方面，树模型表现更稳健，其中 XGBoost 的 AUC 最高，为 0.945；随机森林的准确率与 F1 值最好，分别为 84.9\% 与 0.840。该基准比较表明，在大样本异质表格数据上，深度表格模型和树模型都具有较强竞争力，但树模型在分类任务中更稳定、训练成本更低。

\begin{table}[htbp]
  \centering
  \caption{多类型纳米颗粒基准数据集上的分类性能比较}
  \label{tab:benchmark_classification}
  \begin{tabular}{lcccc}
    \hline
    模型 & Accuracy & F1 & AUC \\
    \hline
    随机森林 & 0.849 & 0.840 & 0.943 \\
    XGBoost & 0.849 & 0.836 & 0.945 \\
    TabTransformer & 0.845 & 0.835 & 0.939 \\
    ABRF & 0.828 & 0.817 & 0.921 \\
    LARF & 0.819 & 0.809 & 0.914 \\
    蒸馏 TabTransformer & 0.722 & 0.707 & 0.801 \\
    \hline
  \end{tabular}
\end{table}

\paragraph{硒纳米颗粒专用数据集上的模型比较。}
在硒纳米颗粒专用数据集上，我们采用贴近实际应用的三折交叉验证方案，并以基础特征与增强特征两种设置分别训练随机森林、XGBoost 与 TabTransformer。基础特征包括浓度、暴露时间、粒径、Zeta 电位、尺寸比、表面电荷、细胞类型、修饰基团和检测方法等；增强特征在此基础上进一步引入 IC50、元素 Se 含量、长期稳定性、Se 3d 结合能 以及晶型编码。考虑到 IC50 与细胞毒性终点之间存在较强关联，以下结果主要用于比较不同增强信息的统计贡献，其中含 IC50 的结果需结合其解释边界加以理解。

结果如表~\ref{tab:senp_main_results} 所示。基础特征下，随机森林与 XGBoost 已具备较好的回归能力，分别达到 $R^2=0.701$ 与 $R^2=0.751$；TabTransformer 在该中小样本场景下明显欠拟合，$R^2$ 仅为 0.003。引入增强变量后，树模型的拟合能力进一步改善：随机森林的三折 CV MAE 从 11.04\% 降至 7.63\%，$R^2$ 从 0.701 提升至 0.915；XGBoost 的三折 CV MAE 从 9.87\% 降至 7.20\%，$R^2$ 进一步提升至 0.957，为所有方案中数值最高。相比之下，TabTransformer 即使加入增强特征后也仅达到 $R^2=0.057$，说明在当前样本规模下，树模型对硒纳米颗粒表格数据更具优势。

方法部分所述的多模态融合策略，即特征拼接、交叉注意力加权与残差连接，应用于全部增强特征配置。在当前数据集上，注意力机制相较简单拼接未带来明显提升，这可能与样本量有限有关；但该机制为引入可选 IC50 描述符及动态调整联合特征表示提供了明确路径。该设计将增强特征中的材料本征化学信息与通用理化描述符整合，对应于上述增强配置的性能改善。

需要强调的是，上述增强设置包含 IC50，因此其性能提升既反映了材料化学信息的补充作用，也可能部分受益于 IC50 与毒性终点之间的天然相关性，因而更适合作为辅助分析结果，而非直接视为真实部署场景下的最终预测性能。

\begin{table}[htbp]
  \centering
  \caption{硒纳米颗粒数据集上基础特征与增强特征的模型性能}
  \label{tab:senp_main_results}
  \begin{tabular}{lccc}
    \hline
    模型 & 特征集 & 3折 CV MAE (\%) & $R^2$ \\
    \hline
    随机森林 & 基础 & 11.04 & 0.701 \\
    随机森林 & 增强 & 7.63 & 0.915 \\
    XGBoost & 基础 & 9.87 & 0.751 \\
    XGBoost & 增强 & 7.20 & 0.957 \\
    TabTransformer & 基础 & 24.13 & 0.003 \\
    TabTransformer & 增强 & 25.50 & 0.057 \\
    \hline
  \end{tabular}
\end{table}

\paragraph{特征效应。}
增强特征模型的特征重要性分析表明，IC50、元素 Se 含量、长期稳定性和Se 3d 结合能 对硒纳米颗粒细胞毒性预测结果具有较大贡献。其中，IC50 直接反映剂量--效应关系，因此更应被视为毒理响应相关辅助变量；元素 Se 含量与颗粒纯度及潜在离子释放背景相关，长期稳定性表征颗粒在分散体系中的聚集与氧化趋势，Se 3d 结合能 则反映表面化学态差异。由此可见，在单一材料体系的毒性建模中，真正具有先验解释意义的增益主要来自能够刻画材料本征化学属性的变量，而不能仅依据含 IC50 模型的高分数判断模型已具备充分的真实预测能力。

\paragraph{关键特征对纳米硒材料选择的指导意义。}

\paragraph{组合特征分析。}
为了识别真正有效的增强信息，我们对五个候选增强变量进行了自动化筛选，候选变量包括 IC50、元素 Se 含量（wt\%）、长期稳定性（周）、Se 3d 结合能（eV）与晶型编码。穷举 31 种非空组合并以三折交叉验证 MAE 为目标后，表现最优的组合为
\[
\{\text{IC50},\ \text{元素 Se 含量},\ \text{长期稳定性},\ \text{Se 3d 结合能}\}.
\]
%% \[
%% \{\mathrm{IC50},\ \mathrm{元素\ Se\ 含量},\ \mathrm{长期稳定性},\ \mathrm{结合能\ Se3d}\}.
%% \]
该组合的三折 CV MAE 为 7.60\%，与使用全部增强特征时的 7.63\% 基本一致，说明晶型编码在当前数据集中的独立贡献有限，而 IC50、Se 含量、稳定性与结合能提供了主要统计增益。需要进一步指出的是，由于该组合包含 IC50，因此其结果更适合作为“含响应相关辅助变量时模型可达到的上界参考”，而不宜简单等同于仅依赖先验理化参数即可实现的预测效果。

\paragraph{模型结果对真实纳米硒材料设计与筛选的指导意义。}

\paragraph{回归与分类任务的最佳模型选择。}
综合上述结果，在硒纳米颗粒专用数据集上，\textbf{XGBoost 增强特征模型}在回归任务中取得了最高数值表现（$R^2=0.957$，CV MAE=7.20\%），高于随机森林和 TabTransformer。在分类任务中，XGBoost 和随机森林在基准数据集上表现相当，XGBoost 的 AUC 略高（0.945 vs 0.943）。但鉴于该回归最优结果对应的增强设置包含 IC50，本文更倾向于将其视为机理增强和辅助分析模型，而非直接作为真实应用场景中的唯一核心预测工具。相较之下，不依赖响应相关辅助变量的基础特征模型及其材料化学增强版本，更适合作为后续外部验证和部署研究的主要比较基线。

\paragraph{与现有工作的比较。}
与 Masarkar 等人 \cite{masarkar2024} 的通用纳米颗粒毒性预测模型相比，本研究具有以下改进：
\begin{itemize}
  \item \textbf{针对单一材料（纳米硒）} 构建了专用数据集，并引入了元素 Se 含量、长期稳定性和Se 3d 结合能 等材料化学增强信息，同时比较了含 IC50 与不含 IC50 的设置差异。这说明单一材料体系中的材料本征化学属性具有重要建模价值，而含 IC50 的高分数结果则更适合作为辅助参考，而非直接与仅依赖先验理化参数的通用模型进行等价比较。
  \item \textbf{系统比较了深度表格模型与树模型}，发现 TabTransformer 在小样本场景（383 条）下严重欠拟合，而树模型（尤其是 XGBoost）表现稳健。这为今后小样本纳米毒性数据集建模提供了参考。
  \item \textbf{通过穷举搜索确定了最优增强特征子集}，剔除了晶型编码等冗余特征，避免了过拟合，同时增强了模型的可解释性。
\end{itemize}
与近期其他基于注意力机制的树模型（如 ABRF、LARF）相比，本研究未发现注意力机制在小规模表格数据上带来显著提升，反而增加了计算开销。因此，对于纳米毒性这类中等规模、混合型特征的数据，传统的梯度提升树（XGBoost）仍是性价比最高的选择。

\paragraph{结论。}
本研究建立的 XGBoost 增强特征模型在当前数据集上表现出较强的拟合能力（$R^2=0.957$），并可同时提供分类决策（高/低存活）。但由于该结果对应的增强设置包含 IC50，故更适合作为辅助分析和机理增强参考。相对而言，元素 Se 含量、长期稳定性与Se 3d 结合能 等材料化学特征所体现的增益，更能支持单一纳米材料毒性建模中的特征工程选择，并为后续外部验证与模型部署提供比较基础。

上述多模态融合框架为整合异质特征模态提供了明确路径；后续将利用独立实验数据开展外部验证，进一步评价其泛化能力。
